\documentclass[12pt]{article}
\usepackage[margin=2.5cm]{geometry}
\usepackage{amsmath,amssymb,amsthm}
\usepackage{bm}

\title{GRA Multiverse Meta-Reset: Formal Definition}
\author{AAA AAA}
\date{March 2026}

\begin{document}
\maketitle

\section{Overview}

This document formalizes the \emph{GRA Multiverse Meta-Reset} (Мета-обнулёнка)
architecture as used in the \texttt{gra-multiverse-asi-lab} repository.
It extends a base GRA reset mechanism to a multi-level, multiverse setting,
and defines a global functional $J_{\text{multiverse}}$ together with
conditions for complete multiverse reset (``обнуление'').[file:850][web:963]

We assume a family of subsystems indexed by multi-indices and organized
into levels $l = 0,1,\dots,K$, where:
\begin{itemize}
  \item Level $0$ contains local domains (tasks, modules, agents).
  \item Level $1$ contains meta-systems over level-$0$ domains.
  \item Level $2$ contains meta-meta-systems, etc.
  \item Level $K$ is the global multiverse level.
\end{itemize}[file:850]

\section{Indexing and State Spaces}

\subsection{Multi-indices}

Let $\mathcal{I}$ be a set of multi-indices
\[
\bm{a} = (a_0, a_1, \dots, a_k),
\]
where $k = \dim(\bm{a})$ denotes the level of $\bm{a}$ in the hierarchy
($0 \le k \le K$).[file:850]

We write
\[
\mathcal{I}^{(l)} = \{ \bm{a} \in \mathcal{I} \mid \dim(\bm{a}) = l \}
\]
for the set of indices at level $l$.

We use the relation
\[
\bm{b} \prec \bm{a}
\]
to mean that $\bm{b}$ is a \emph{subsystem} of $\bm{a}$; formally, this specifies
the hierarchical embedding of lower-level subsystems into higher-level ones.[file:850]

\subsection{Local Hilbert Spaces and Level Spaces}

To each index $\bm{a} \in \mathcal{I}$ we associate a Hilbert space
$\mathcal{H}^{(\bm{a})}$, representing the state space of subsystem $\bm{a}$.

For each level $l$ we define the level space
\[
\mathcal{H}^{(l)} =
\bigotimes_{\bm{a} \in \mathcal{I}^{(l)}} \mathcal{H}^{(\bm{a})}.
\]

The total multiverse Hilbert space is
\[
\mathcal{H}_{\text{multiverse}} =
\bigotimes_{l=0}^K \mathcal{H}^{(l)}.
\]

A multiverse state is a vector
\[
\Psi_{\text{multiverse}} \in \mathcal{H}_{\text{multiverse}},
\]
which can be decomposed into its level components and subsystem components
\[
\Psi_{\text{multiverse}}
\longleftrightarrow
\{ \Psi^{(\bm{a})} \}_{\bm{a} \in \mathcal{I}}.
\][file:850]

\section{Goals, Projectors, and Foam}

\subsection{Goals and Projectors at Each Level}

For each level $l$ we define:
\begin{itemize}
  \item A goal $G_l$ (or a family $\{G_l^{(\bm{a})}\}_{\bm{a} \in \mathcal{I}^{(l)}}$).
  \item A projector $\mathcal{P}_{G_l}$ (or $\mathcal{P}_{G_l^{(\bm{a})}}$) onto the
        solution subspace corresponding to that goal.
\end{itemize}

These projectors encode what it means for a subsystem or collection of subsystems
to be in a ``goal-satisfying'' state at level $l$.[file:850]

\subsection{Foam at Level \texorpdfstring{$l$}{l}}

Let $\Psi^{(l)}$ denote the joint state at level $l$, interpreted as the tensor
product of $\Psi^{(\bm{a})}$ over all $\bm{a} \in \mathcal{I}^{(l)}$.

We define the \emph{foam} at level $l$ as
\[
\Phi^{(l)}\big(\Psi^{(l)}, G_l\big) =
\sum_{\bm{a} \neq \bm{b} \atop \bm{a},\bm{b} \in \mathcal{I}^{(l)}}
\left|
  \left\langle
    \Psi^{(\bm{a})}
    \,\middle|\,
    \mathcal{P}_{G_l}
    \,\middle|\,
    \Psi^{(\bm{b})}
  \right\rangle
\right|^2.
\]

Intuitively, $\Phi^{(l)}$ measures off-diagonal, goal-unexplained correlations
between subsystems at level $l$; non-zero values indicate interpretational
artifacts or inconsistencies.[file:850]

\section{Multiverse Functional}

\subsection{Local Functional at Level 0}

At level $0$, for each $\bm{a} \in \mathcal{I}^{(0)}$ we define a local
functional
\[
J^{(0)}\big(\Psi^{(\bm{a})}\big)
=
J_{\text{loc}}\big(\Psi^{(\bm{a})}; G_0^{(\bm{a})}\big),
\]
which encodes task performance, base alignment, and other local criteria.[file:850]

\subsection{Recursive Definition of Level Functionals}

For $l \ge 1$ and $\bm{a} \in \mathcal{I}^{(l)}$ we define recursively:
\[
J^{(l)}\big(\Psi^{(\bm{a})}\big)
=
\sum_{\substack{\bm{b} \prec \bm{a} \\
                \bm{b} \in \mathcal{I}^{(l-1)}}}
J^{(l-1)}\big(\Psi^{(\bm{b})}\big)
+
\Phi^{(l)}\big(\Psi^{(\bm{a})}, G_l^{(\bm{a})}\big).
\]

Here the sum over $\bm{b} \prec \bm{a}$ aggregates contributions from all
subsystems contained in $\bm{a}$, and the foam term penalizes inconsistency
at level $l$.[file:850]

\subsection{Global Multiverse Functional}

We introduce non-negative weights for levels
\[
\Lambda_l = \lambda_0 \alpha^l,
\quad 0 < \alpha < 1,
\]
with a base weight $\lambda_0 > 0$.[file:850]

The full multiverse functional is
\[
J_{\text{multiverse}}(\bm{\Psi}) =
\sum_{l=0}^K \Lambda_l
\sum_{\bm{a} \in \mathcal{I}^{(l)}}
J^{(l)}\big(\Psi^{(\bm{a})}\big),
\]
where $\bm{\Psi} = \{\Psi^{(\bm{a})}\}_{\bm{a} \in \mathcal{I}}$ is the collection
of all subsystem states.

Minimizing $J_{\text{multiverse}}$ corresponds to simultaneously:
\begin{itemize}
  \item solving local tasks,
  \item aligning subsystems and meta-systems,
  \item eliminating hierarchical foam across all levels.[web:963][web:956]
\end{itemize}

\section{Complete Multiverse Reset}

\subsection{Definition}

We say that the multiverse is in a state of \emph{complete reset}
if there exists $\bm{\Psi}^* = \{\Psi^{(\bm{a})*}\}_{\bm{a}}$ such that
\[
\Phi^{(l)}\big(\Psi^{(l)*}, G_l\big) = 0
\quad \forall\, l = 0,\dots,K.
\]

Equivalently, at each level $l$ all off-diagonal, goal-unexplained
correlations vanish, and each subsystem is perfectly aligned
with its goal projector.[file:850]

\subsection{Assumptions}

We consider the following structural conditions:
\begin{enumerate}
  \item \textbf{Commutativity of goal projectors}:
    \[
    \big[
      \mathcal{P}_{G_l^{(\bm{a})}},
      \mathcal{P}_{G_m^{(\bm{b})}}
    \big] = 0
    \quad
    \forall\, \bm{a},\bm{b},\ l,m.
    \]
  \item \textbf{Hierarchical consistency}:
    for $l \ge 1$ and $\bm{a} \in \mathcal{I}^{(l)}$,
    \[
    \mathcal{P}_{G_l^{(\bm{a})}} =
    \bigotimes_{\substack{\bm{b} \prec \bm{a} \\
                          \bm{b} \in \mathcal{I}^{(l-1)}}}
    \mathcal{P}_{G_{l-1}^{(\bm{b})}}.
    \]
  \item \textbf{Sufficient dimension}:
    \[
    \dim(\mathcal{H}_{\text{multiverse}})
    \ge
    \prod_{l=0}^K N_l,
    \]
    where $N_l = |\mathcal{I}^{(l)}|$ is the number of subsystems at level $l$.[file:850]
\end{enumerate}

\subsection{Theorem (Multiverse Reset)}

Under the above conditions there exists a state
$\bm{\Psi}^* = \{\Psi^{(\bm{a})*}\}$ such that
\[
\Phi^{(l)}\big(\Psi^{(l)*}, G_l\big) = 0
\quad \forall\, l = 0,\dots,K.
\]

\paragraph{Proof sketch.}
\begin{itemize}
  \item \textbf{Base step ($l=0$):} follows from the local reset theorem
        (existence of states aligned with $\mathcal{P}_{G_0^{(\bm{a})}}$ for all $\bm{a} \in \mathcal{I}^{(0)}$).[file:850]
  \item \textbf{Inductive step:} assume reset achieved up to level $l-1$.
    Using hierarchical consistency, we have
    \[
    \mathcal{P}_{G_l^{(\bm{a})}} =
    \bigotimes_{\bm{b} \prec \bm{a}}
    \mathcal{P}_{G_{l-1}^{(\bm{b})}}.
    \]
    If each $\Psi^{(\bm{b})*}$ is an eigenvector of $\mathcal{P}_{G_{l-1}^{(\bm{b})}}$,
    then their tensor product is an eigenvector of $\mathcal{P}_{G_l^{(\bm{a})}}$.
    Working in a basis where all projectors are diagonal (commutativity),
    off-diagonal terms vanish and thus $\Phi^{(l)} = 0$.
\end{itemize}

\section{Optimization and Dynamics}

\subsection{Gradient Structure}

The gradient of $J_{\text{multiverse}}$ with respect to subsystem state
$\Psi^{(\bm{a})}$ at level $l = \dim(\bm{a})$ can be written as
\[
\frac{\partial J_{\text{multiverse}}}{\partial \Psi^{(\bm{a})}} =
\Lambda_l
\frac{\partial \Phi^{(l)}}{\partial \Psi^{(\bm{a})}}
+
\sum_{\bm{b} \succ \bm{a}}
\Lambda_{\dim(\bm{b})}
\frac{\partial \Phi^{(\dim(\bm{b}))}}{\partial \Psi^{(\bm{a})}},
\]
where the sum is over all higher-level subsystems $\bm{b}$ that contain $\bm{a}$.[file:850]

A basic gradient descent update has the form
\[
\Psi^{(\bm{a})}(t+1) =
\Psi^{(\bm{a})}(t)
-
\eta
\left[
\Lambda_l
\frac{\partial \Phi^{(l)}}{\partial \Psi^{(\bm{a})}}
+
\sum_{\bm{b} \succ \bm{a}}
\Lambda_{\dim(\bm{b})}
\frac{\partial \Phi^{(\dim(\bm{b}))}}{\partial \Psi^{(\bm{a})}}
\right],
\]
with learning rate $\eta > 0$.

\subsection{Complexity}

For a multiverse with $K$ levels and $N$ subsystems per level (for simplicity),
the total cost of computing pairwise foam terms is
\[
\mathcal{O}\left(
  \sum_{l=0}^K N^2 \alpha^l
\right)
=
\mathcal{O}\left(
  N^2 \frac{1 - \alpha^{K+1}}{1 - \alpha}
\right),
\quad 0 < \alpha < 1.
\]

In the limit $K \to \infty$ and fixed $\alpha < 1$, this remains polynomial:
\[
\mathcal{O}\left(
  \frac{N^2}{1-\alpha}
\right).
\][file:850]

\section{Limit State and Cognitive Vacuum}

Consider the limit $K \to \infty$ (infinite hierarchy) and a sequence of optimal
states $\Psi_K^*$ for each finite $K$. If the limit
\[
\Psi_\infty^* = \lim_{K \to \infty} \Psi_K^*
\]
exists in an appropriate sense, and
\[
\Phi^{(l)}\big(\Psi_\infty^*, G_l\big) = 0
\quad \forall\, l \in \mathbb{N},
\]
then $\Psi_\infty^*$ lies in the intersection
\[
\Psi_\infty^* \in \bigcap_{l=0}^\infty \ker(\Phi^{(l)}).
\]

This intersection can be interpreted as a \emph{cognitive vacuum}:
a state free from interpretational artifacts at all levels, containing
only structural invariants of the underlying reality.[file:850][web:963]

\section{Remarks on Safety and Self-Modification}

In the context of self-improving agents, we restrict allowed
modifications $m$ of the state and architecture to a safe set
\[
\mathcal{M}_{\text{safe}} =
\{ m \mid [\mathcal{P}_{G_l}, m] = 0 \ \forall l \},
\]
ensuring that goal projectors remain fixed while the agent optimizes
its internal structure with respect to $J_{\text{multiverse}}$.[web:961][web:964][web:956]

This makes the multiverse reset functional suitable as a target for
recursive self-improvement without trivial goal corruption or wireheading.

\end{document}
